\section{From validated readouts to complete deployment}\label{sec:deployment70}
The fixed-feature construction and its optimization certificate identify an achieved numerical object. The next issue is its use after training. Independent validation need not imply that the states subsequently visited by a controller are independent. We give a work account that keeps those two sources of information separate and charges the descendants of a failed screen.

\subsection{Predictable selection and changing context laws}
Condition on a frozen family $\Theta$ of $M$ readouts. Write
\[
 f_\theta(z)=K(z)\ell_\tau(z,\theta),\qquad
 0\le N_\theta(z)\le f_\theta(z)\le K_{\max},
\]
where the fallback and its endpoint cap are the ones specified for the context. Let $Z_1^v,\ldots,Z_n^v$ be independent validation contexts with law $P$, independent of readout construction. Set
\begin{equation}\label{eq:validated-upper70}
 u_\theta=\min\left\{K_{\max},\ \frac1n\sum_{i=1}^nf_\theta(Z_i^v)
       +K_{\max}\sqrt{\frac{\log(M/\alpha)}{2n}}\right\}.
\end{equation}
During deployment, let $\mathcal F_{b-1}$ contain construction, validation, and all observations before block $b$. The selected readout $\theta_b$ is $\mathcal F_{b-1}$-measurable. The new context $Z_b$ can depend on this history. Suppose a declared predictable number $\rho_b\in[0,1]$ satisfies
\begin{equation}\label{eq:context-transfer70}
 \|\mathcal L(Z_b\mid\mathcal F_{b-1})-P\|_{\rm TV}\le\rho_b,
 \qquad \|P-Q\|_{\rm TV}=\sup_A|P(A)-Q(A)|.
\end{equation}
This is a distributional condition, not a consequence of matching a model hash.

Use an additive operation model. Let $F_0$ bound fitting, serialization and validation work. Let $h_b$ bound loading, identity checks, prediction, previously purchased endpoints, acquisition and durable output at block $b$. Finally let a nonrandom $c_b$ bound the complete descendant and numerical-recovery work of one additional root endpoint. The budgets must cover all admissible histories, not merely the largest observed runtime. Assume the deployed work satisfies
\[
 W_B\le F_0+\sum_{b=1}^B\{h_b+c_bN_{\theta_b}(Z_b)\}.
\]
This includes a complete-action mesh fallback when that is the fallback validated by the experiment. Substitution of a different fallback requires its own cap and loss contract.

\begin{theorem}[Validated work under predictable deployment]\label{thm:deployment70}
For a fixed deployment window $B$ and $0<\alpha,\gamma<1$, with probability at least $1-\alpha-\gamma$,
\begin{align}\label{eq:deployment-work70}
 W_B\le {}&F_0+\sum_{b=1}^B h_b
 +\sum_{b=1}^Bc_b\min\{K_{\max},u_{\theta_b}+K_{\max}\rho_b\}\notag\\
 &+K_{\max}\sqrt{\frac12\log(1/\gamma)\sum_{b=1}^Bc_b^2}.
\end{align}
The readout can be selected using the validation sample and the preceding deployment history. The conclusion does not require the deployment contexts to be independent. Under the same local feasibility and certificate premises, this selection does not change the original-optimum policy guarantee.
\end{theorem}
\begin{proof}
The bounded-sum argument in Theorem~\ref{thm:work-loss67} gives $P f_\theta\le u_\theta$ for every $\theta$ on one event of probability at least $1-\alpha$. For $0\le f\le K_{\max}$, the layer-cake representation gives $|Pf-Qf|\le K_{\max}\|P-Q\|_{\rm TV}$. Hence, on that event,
\[
 \mathbb E[N_{\theta_b}(Z_b)\mid\mathcal F_{b-1}]
 \le\min\{K_{\max},u_{\theta_b}+K_{\max}\rho_b\}.
\]
The variables $c_bN_{\theta_b}(Z_b)$ have conditional ranges of length at most $c_bK_{\max}$. Conditional Hoeffding bounds, multiplied successively along the filtration and followed by the exponential Markov inequality, show that the sum exceeds its conditional means by at most the final term of \eqref{eq:deployment-work70}, except on an event of probability $\gamma$. Union with the validation failure event proves the result. Every selected readout remains a proposal to the same verified gate. Its statewise certificate, and thus the earlier acquired-policy recursion, is unchanged.
\end{proof}

This theorem gives a prospective route from training to complete work rather than interpreting a training loss as a timing result. A smaller root-query mean is useful when its descendant-weighted savings exceed setup and per-block overhead. The $c_b$ terms expose the continuation-tree cost. Setting them to one would prove a root-query statement, not a full-service statement. The R69 independent-context experiment below uses its declared mesh recovery; its adaptive recursive deployment records are not claimed to satisfy an empirically verified version of \eqref{eq:context-transfer70}.

\subsection{Policy identity, declined reuse, and adoption}
Cost inference can be combined with a deployment clock only after identifying the controller whose cost was estimated. The release checks the complete readout arrays, tolerance, acquisition precision, action rounding and original kernels. Consider a service processing $B$ blocks of $m$ trajectories each. In $D$ predetermined blocks it declines an incompatible readout and uses the benchmark $\nu$; otherwise it uses a fixed $\pi$. Suppose each block has the same initial law and economic primitives, and the decline schedule is independent of the trajectories. Let $[l,u]$ be a simultaneous interval for $\Delta=\mathbb E[J^\pi-J^\nu]$.

\begin{proposition}[Complete-service adoption region]\label{prop:adoption70}
Let $W_\pi,W_\nu$ be complete work observations under the same boundary. If $v\ge0$ values one normalized expected-cost unit, $p\ge0$ prices one work unit, and $f\ge0$ is the incremental installation fee, then the net cost difference belongs to
\begin{equation}\label{eq:adoption-region70}
 [\,vm(B-D)l+p(W_\pi-W_\nu)+f,
    \ vm(B-D)u+p(W_\pi-W_\nu)+f\,]
\end{equation}
on the original simultaneous event. The statement holds for every such price triple on that event. With observed work saving $S=W_\nu-W_\pi>0$, strict adoption is supported whenever
\[
 pS>f+vm(B-D)u.
\]
If work is instead bounded by valid intervals, replacing $W_\pi-W_\nu$ by its interval difference preserves the statement.
\end{proposition}
\begin{proof}
The expected cost difference in a normal block is $m\Delta$; it is zero in a declined block. Linearity of expectation gives $m(B-D)\Delta$, without requiring block independence. Multiply the cost interval by its nonnegative price, and add the work difference and fee. All price choices use the same event, so no additional probability allocation is needed.
\end{proof}

The executed reuse design has $B=64$, $m=8$, and one predetermined identity rejection, so its normal-policy multiplier is $504$. This does not make a observed prefix clock a completed service, or an observed timing range a confidence interval for future hardware. The paired complete receipts and the conservative envelope across the observed repetitions are both reported. The inputs $v,p,f$ remain declared decision parameters; the experiment estimates neither an installation fee nor the monetary value of normalized investment losses. The policy-cost margin supplies an economically interpretable eligibility condition, while the full adoption region keeps computational savings and policy quality separate.
